\scanpage{219}
\boxed{\text{Quelques complexes d'algèbres différentielles.}}

$A$ anneau commutatif, $M$ $A$-module.
\margin{On considère $\operatorname{Sym}^*(M)\otimes\Lambda^*M$ et $\Gamma^*(M)\otimes\Lambda^*M$ \uline{comme} $A$-alg graduées par \uline{degré de} $\Lambda M$}
\margin{De Rham}
\textcircled{1} Sur $\operatorname{Sym}^*(M)\otimes\Lambda^*M$ \deleted{(algèbre bigraduée)}, il y a un unique opérateur $d$ (qui soit \deleted{a) $A$ linéaire} a) une \inserted{anti-}\uline{dérivation} \inserted{de $A$-algèbre graduée} \deleted{(pour $\operatorname{Sym}^*M$ de degré pair et \deleted{$\operatorname{Sym}$} $\Lambda M$ de degré impair)} \inserted{b)} qui soit nulle sur $\Lambda^1M$, et qui sur $\operatorname{Sym}^1(M)$ soit l'isom can. $\operatorname{Sym}^1(M)\longrightarrow\Lambda^1M$.

Ce $d$ \deleted{\unsure{est}} satisfait $d^2=0$, et applique $\operatorname{Sym}^p\otimes\Lambda^q$ dans $\operatorname{Sym}^{p-1}\otimes\Lambda^{q+1}$. \inserted{Il est $\Lambda^*M$-linéaire.} \deleted{\unsure{Est $A$ linéaire}} C'est aussi le complexe de De Rham de $\operatorname{Sym}^*(M)$ sur $A$, donc le complexe de De Rham du fibré vectoriel $V(M)$ sur $A$. \deleted{Si $\deleted{\unsure{M}}M$ est libre et Si $M$ est libre},
Supposons $M$ \unsure{plat}, \deleted{(plus gén., plat)} $A$ de car. $0$ (i.e. $A$ une $\mathbf Q$-algèbre), alors
\margin{inutile}
\deleted{$H^*(\operatorname{Sym}^*(M)\otimes\Lambda^*M)$ (gradué par le degré de $\Lambda^*(M)$) est nul en degrés \unsure{$\neq0$ et en degré $0$}}
(graduant suivant le degré de $\Lambda^*M$
\[
H^i_{DR}(\operatorname{Sym}^*(M)\otimes\Lambda^*(M))\simeq
\begin{cases}0&\text{si }i\neq0\\
\underset{\substack{\Vert\\\operatorname{Sym}^0\otimes\Lambda^0}}{A}&\text{si }i=0
\end{cases}
\]
Si $A$ est de car. $p>0$, \inserted{(i.e. est une $\mathbf F_p$-alg), \unsure{$A$ plat|$M$ plat}} alors on a un (histoire des opérations, de Cartier ----)

\margin{Koszul}
\textcircled{2} Sur $\operatorname{Sym}^*(M)\otimes\Lambda^*M$ \deleted{(algèbre bigraduée)} il y a un unique opérateur $\delta$ qui est \deleted{a) $A$ linéaire} a) une \inserted{anti-}dérivation \deleted{qui} de $A$-alg. \unsure{graduées|graduée} et \deleted{\unsure{resp. pour} $\operatorname{Sym}^1M$ de degrés pairs et b) $\Lambda^*M$ de degrés impairs}, b) nulle sur \deleted{$\unsure{A}$} $\operatorname{Sym}^1M$, et sur $\Lambda^1M$ c'est l'isom can. $\Lambda^1M\xrightarrow{\sim}\operatorname{Sym}^1M$. Alors $\delta$ satisfait $\delta^2=0$, il applique $\operatorname{Sym}^p\otimes\Lambda^q$ dans $\operatorname{Sym}^{p+1}\otimes\Lambda^{q-1}$. Si $M$ est libre (plus gén., plat) et sans condition de car. sur $A$, on a \deleted{\unsure{que}} \inserted{(graduant} la ch. par le degré de $\Lambda^*$ \inserted{qu'on écrit en bas} (de sorte que $\delta$ est de degré \deleted{$1$}\inserted{$\unsure{i-1.1|-1}$}) :
\[
H^{\mathrm{Kos}}_{\unsure{q|i}}(\operatorname{Sym}^*M\otimes\Lambda^*(M))=
\begin{cases}0&\text{si }i\neq0\\
\simeq\underset{\substack{\Vert\\\operatorname{Sym}^0\otimes\Lambda^0}}{A}&\text{si }i=0
\end{cases}.
\]
\margin{218}
